\subsection{Schopnost pojmout celé studentské řešení}
\label{subsec:criteria-context-window}

V \secref[Sekci]{subsec:llm-token-limit} byla popsána omezená kapacita kontextového okna LLM, která určuje, kolik tokenů může model zpracovat v jednom požadavku.
Pro analýzu studentských řešení je klíčové, aby model dokázal pojmout celé odevzdání včetně zadání a zdrojového kódu, protože pouze tak může poskytnout relevantní a kontextově správnou zpětnou vazbu.
Zároveň mu musí zůstat dostatek prostoru pro generování kvalitních komentářů.
Je tedy nutné pojmout následující části:

\begin{enumerate}
    \item \textbf{Systémový prompt} -- obsahuje instrukce pro model, které definují jeho roli, požadavky na formát výstupu a další klíčové parametry generování.
    Typicky se jedná o několik strukturovaných vět s jasnými pokyny, které mohou zabrat 200--500 tokenů.

    \item \textbf{Zadání úlohy} -- popis úlohy, kterou student řeší, včetně požadavků na vstupy, výstupy a případné specifické podmínky.
    Zadání může být poměrně rozsáhlé, zejména pokud obsahuje příklady vstupů a výstupů, což může zabrat 500--1~000 tokenů.

    \item \textbf{Zdrojový kód} -- samotné řešení, které student odevzdal.
    Velikost zdrojového kódu se může velmi lišit, ale pro středně velkou úlohu může být v rozsahu 1~000--5~000 tokenů.
\end{enumerate} \TODO{Ověřit konkrétní hodnoty}

Pro standardní studentské odevzdání v systému Kelvin je tedy dostačující kontextové okno přibližně \textbf{8~000 tokenů}.
Aby byl systém schopen zvládnout i rozsáhlejší úlohy (větší projekty, vícesouborová řešení) a aby byl respektován princip dostatečné rezervy, je jako \textbf{minimální požadavek} stanoveno \textbf{16~000 tokenů}.
Moderní modely tuto hranici zpravidla bez problémů splňují, většina aktuálně dostupných modelů nabízí kontextové okno 32~000 až 128~000 tokenů.

\begin{table}[H]
    \centering
    \resizebox{\textwidth}{!}{%
        \begin{tabular}{lp{9cm}l}
            \toprule
            \textbf{Kritérium}              & \textbf{Popis}                                                                 & \textbf{Minimální požadavek} \\
            \midrule
            Kvalita výstupů                 & Věcně správné, srozumitelné, relevantní komentáře                              & Vyšší než průměr \\
            Dostupnost                      & Přístupný přes API nebo ke stažení; vhodná licence pro akademické použití      & Open-source nebo API s DPA \\
            Náročnost provozu               & Přijatelné náklady na token nebo na hardware                                   & Viz \tabref{tab:deployment-comparison} \\
            Kontextové okno                 & Schopnost zpracovat celé odevzdání včetně promptu                              & $\geq$ 16~000 tokenů \\
            \bottomrule
        \end{tabular}
    }
    \caption{Přehled kritérií výběru modelu a jejich minimálních požadavků}
    \label{tab:selection-criteria}
\end{table}

\endinput